\begin{table}[!htbp]
    \centering
    \caption{Kriteria penilaian alternatif pendekatan beserta bobot dan asal kebutuhan.}
    \label{tab:tab3.4_kriteria_penilaian}
    \renewcommand{\arraystretch}{1.15}
    \begin{tabular}{@{}p{1.0cm}p{9.7cm}p{1.4cm}@{}}
        \hline
        \textbf{Kode} & \textbf{Kriteria dan asal kebutuhan} & \textbf{Bobot} \\
        \hline
        C1 & Kemampuan antisipatif, yaitu menyediakan taksiran kuantitatif kondisi yang akan datang beserta ketidakpastiannya (KF-03, KF-07). & 0,25 \\
        C2 & Pengondisian rekomendasi pada kondisi terprediksi, yaitu rekomendasi disusun dari kondisi yang diperkirakan akan terjadi, bukan dari kondisi yang sedang berlaku (KF-04, KF-06). & 0,25 \\
        C3 & Keterlacakan dan keterjelasan, yaitu keluaran dapat ditelusuri sampai ke dokumen dan halaman sumbernya (KNF-02, KNF-08). & 0,20 \\
        C4 & Keterterapan tanpa infrastruktur berat, yaitu berjalan tanpa akselerator grafis, tanpa sensor pemantauan sinambung, dan tanpa server khusus (KNF-01, KNF-09). & 0,20 \\
        C5 & Keamanan klinis dan keamanan konten, yaitu keluaran dibumikan pada dokumen dan dokter tetap menjadi penentu akhir (KNF-03, KNF-06). & 0,10 \\
        Total & & 1,00 \\
        \hline
    \end{tabular}
\end{table}
